\section{\bt{Introduction}{引言}}
\label{sec:introduction}
\bi{Large language models (LLMs) are enabling data agents for database analysis and predictive modeling~\cite{lei2025spider2,guo2024dsagent,liu2026agents}. These tasks rely on database-specific measures, grain, filters, and join paths~\cite{floratou2024nl2sql,robinson2024relbench}. For example, March, September, and third-quarter revenue all share one underlying definition (Figure~\ref{fig:metric-record}). When sessions run independently, each repeats schema inspection, value probes, join tests, and the associated LLM calls. Across 100 memoryless metric sessions, 93\% of lookups target facts found by earlier sessions, and 71\% of session time precedes the final answer query. Sharing these discoveries across agents and sessions can amortize this exploration cost.}
{大语言模型（LLM）正推动面向数据库分析与预测建模的数据智能体发展~\cite{lei2025spider2,guo2024dsagent,liu2026agents}。这些任务依赖数据库特有的度量、粒度、过滤条件与连接路径~\cite{floratou2024nl2sql,robinson2024relbench}。例如，3 月、9 月与第三季度的营业额共用同一定义（图~\ref{fig:metric-record}）。会话独立运行时，每个会话都会重复执行模式检查、取值探查、连接测试及相关的大模型调用。在我们追踪的 100 个无记忆指标会话中，93\% 的查找以此前会话已获得的事实为目标，71\% 的会话时间消耗在最终答案查询开始之前。跨智能体和会话共享这些发现，能够分摊这部分探索开销。}

\begin{figure}[t]
\centering
% Vector drawing generated by tools/figures/metric_record.py; rerun it after changing the figure.
\ifbilingual
\includegraphics[width=\columnwidth]{figures/metric-record-zh}
\else
\includegraphics[width=\columnwidth]{figures/metric-record}
\fi
\caption{\bt{Shared memory reduces repeated exploration. (a) Agents independently inspect schemas, probe values, and test joins. (b) Successful sessions publish definitions to \system; later agents reuse them, sharing validation and maintenance. Savings are from the \PpAgents-agent TPC-DS study.}{共享记忆减少重复探索。（a）智能体各自检查模式、探查取值并验证连接。（b）成功会话向 \system\ 发布定义，后续智能体复用定义并共享验证与维护。图中收益来自 \PpAgents\ 个智能体的 TPC-DS 实验。}}
\label{fig:metric-record}
\Description{A before-and-after comparison sharing one database. At left, three agents A, B, and C each have a separate stack of schema, value-probe, and join-test icons connected to the database. At right, the three agents ask March, September, and quarterly revenue questions above a prominent MAVRA shared-memory layer, which sits above the database. A successful session contributes a revenue definition and later agents reuse it for their questions. The layer shows revenue = SUM(paid), grain, filters, and joins; a smaller internal strip contains validation and maintenance. SQL flows between the layer and database, while database updates feed maintenance. The footer reports the population study with \PpAgents\ agents on TPC-DS, relative to no memory: \PpMavraTokSaveNone\ percent fewer tokens and a \PpMavraSpeedupNone-fold end-to-end speedup.}
\end{figure}


\bi{Prior work reduces repeated work through agent memory and explicit database semantics. Agent memories retain insights, workflows, and playbooks~\cite{zhao2024expel,wang2025awm,zhang2026ace}, or share trajectories and composite tools~\cite{zhang2025gmemory,biswal2026agentsm}. Database-oriented systems recover calculation logic from scripts~\cite{weng2025datalab} and centralize metric definitions for SQL generation~\cite{dbt2026metricflow,databricks2026metricviews}. These approaches provide knowledge that later sessions can reuse.}
{既有工作通过智能体记忆与显式数据库语义减少重复工作。智能体记忆保留洞见、工作流与经验手册~\cite{zhao2024expel,wang2025awm,zhang2026ace}，或共享轨迹与复合工具~\cite{zhang2025gmemory,biswal2026agentsm}。面向数据库的系统从脚本恢复计算逻辑~\cite{weng2025datalab}，集中管理指标定义以生成 SQL~\cite{dbt2026metricflow,databricks2026metricviews}。这些方法为后续会话提供了可复用的知识。}

\bi{However, retrieving prior knowledge can still leave each agent to inspect the database and reconstruct SQL for its own question. For the revenue questions above, agents need the same measure, grain, joins, and filters, instantiated for different periods. Sharing this calculation as a parameterized executable definition lets later sessions bypass repeated schema inspection, value probes, join tests, and SQL construction. Its discovery and publication costs can then be spread across all consumers. Data changes introduce further work: the definition's assumptions must be rechecked and the definition maintained. Shared validation and maintenance extend the same amortization to data updates.}
{然而，检索到既有知识后，智能体仍可能需要检查数据库，并为自己的问题重新构造 SQL。对于上述营业额问题，智能体需要相同的度量、粒度、连接与过滤，仅使用的期间不同。将这一计算保存为参数化的可执行定义，便能让后续会话省去重复的模式检查、取值探查、连接测试与 SQL 构造，发现和发布定义的开销也由所有使用者分摊。数据变化会带来进一步的工作：需要重验定义的前提并维护定义。共享验证与维护，将同样的分摊方式延伸到数据更新。}

\bi{Our goal is to reduce the total cost of repeated database exploration across agents, tasks, and sessions. This entails three challenges. \emph{First, extracting reusable computation.} Successful traces mix business logic, such as which sales rows qualify, with task parameters, such as a month or prediction cutoff. Extraction must preserve the former while parameterizing the latter. \emph{Second, making discoveries usable by later agents.} New sessions must find the relevant calculation and instantiate it for their periods or prediction instances, preserving its meaning while avoiding repeated exploration and SQL construction. \emph{Third, sustaining reuse at low cost.} Data updates can invalidate a calculation's assumptions. Revalidating independently for every agent duplicates checks on shared dependencies, while discarding definitions after every write repeats the original exploration. Validation and maintenance must therefore be shared across consumers.}
{我们的目标是降低不同智能体、任务和会话中重复探索数据库的总成本。这面临三大挑战。\emph{第一，提取可复用的计算。}成功轨迹混合了业务逻辑（如哪些销售记录符合统计条件）与任务参数（如月份或预测截点）；提取过程必须保留前者，并将后者参数化。\emph{第二，让后续智能体用得上已有发现。}新会话需要找到相关计算，为自身期间或预测样本实例化，在保留计算含义的同时省去重复探索和 SQL 构造。\emph{第三，以低成本维持复用。}数据更新可能使计算所依赖的前提失效。为每个智能体独立重验，会在共同依赖上重复检查；每次写入后丢弃定义，又会重复最初的探索。因此，验证与维护工作也需要在使用者之间共享。}

\bi{We present \system\ (\textbf{M}etric-\textbf{A}ware \textbf{V}alidation and \textbf{R}euse for \textbf{A}gents), a shared memory layer between data agents and the database that reduces repeated exploration by making session discoveries executable for later agents. Successful sessions contribute metric and feature definitions, table descriptions, and verified join paths to a database-wide library. A definition records the measure, grain, joins, and filters separately from question parameters; feature definitions also record the entity and temporal window. Deterministic publication checks validate extracted definitions and compare compiled results with verified metric answers or submitted feature values. Later agents retrieve accepted definitions and compile SQL for new periods and prediction instances, replacing the exploration that produced the original calculation.}
{本文提出 \system{}（\textbf{M}etric-\textbf{A}ware \textbf{V}alidation and \textbf{R}euse for \textbf{A}gents），一个位于数据智能体与数据库之间的共享记忆层，将会话中的发现转化为后续智能体可执行的知识，减少重复探索。成功会话向数据库范围的共享库贡献指标与特征定义、表描述和经验证的连接路径。定义将度量、粒度、连接和过滤与题目参数分离，特征定义还记录实体与时间窗口。确定性的发布检查验证提取结果，并将编译结果与经过核实的指标答案或提交的特征值比较。后续智能体检索已通过检查的定义，为新的期间与预测样本编译 SQL，省去获得原始计算所需的探索。}

\bi{\system\ amortizes learning and publication costs across consumers of each accepted definition. To sustain reuse after updates, condition identities and dependency versions let agents share validation results, recheck affected conditions, and reuse validated metric repairs. Snapshot-bound metric execution aligns these checks with the consuming query's data. Together, these mechanisms spread discovery, SQL construction, validation, and maintenance over the workload, allowing later agents to reuse earlier work with little additional exploration (Section~\ref{sec:model}).}
{\system\ 将每条已发布定义的学习和发布成本分摊到其使用者。为在更新后维持复用，系统通过条件标识与依赖版本共享验证结果，重验受影响的条件，并复用经过验证的指标修复。快照绑定的指标执行将这些检查与消费查询读取的数据对齐。这些机制将发现、SQL 构造、验证与维护开销分摊到整个负载，使后续智能体只需少量额外探索即可复用此前的工作（第~\ref{sec:model} 节）。}

\bi{We evaluate \system\ on RelBench~\cite{robinson2024relbench}, TPC-DS, and TPC-H. Across \PpRuns\ runs of the TPC-DS workload, \PpAgents\ agents use \PpMavraTokSaveNone\% fewer tokens and run \PpMavraSpeedupNone$\times$ faster than without memory, using \PpMavraTokSavePrivate\% fewer tokens than private memory with the same mechanisms. Post-restatement accuracy reaches \PpMavraChangeAcc\%, compared with \PpNoneChangeAcc\% without memory. In the five-cutoff RelBench study, reuse-stage token consumption (cutoffs 2--5) falls by \RbStackMavraTokSave\% on rel-stack and \RbTrialMavraTokSave\% on rel-trial relative to no memory; all six tasks retain baseline predictive quality within the evaluation margins. Including initial learning and memory operations across all five cutoffs, token use drops by \RbStackMavraTokSaveAll\% and \RbTrialMavraTokSaveAll\%, respectively. A separate three-cutoff study also reduces reuse-stage tokens by \OfStackSaveReuse\% and \OfTrialSaveReuse\%; Section~\ref{sec:relbench-eval} reports both studies under their respective configurations.}
{我们在 RelBench~\cite{robinson2024relbench}、TPC-DS 和 TPC-H 上评估 \system。在 TPC-DS 负载的 \PpRuns\ 轮运行中，\PpAgents\ 个智能体相对无记忆减少 \PpMavraTokSaveNone\% 的 token，快 \PpMavraSpeedupNone\ 倍；相对采用相同机制的私有记忆减少 \PpMavraTokSavePrivate\% 的 token。数据更正后，正确率达到 \PpMavraChangeAcc\%，无记忆组为 \PpNoneChangeAcc\%。在五截点 RelBench 实验中，复用阶段（截点 2--5）的 token 相对无记忆在 rel-stack 与 rel-trial 上分别减少 \RbStackMavraTokSave\% 与 \RbTrialMavraTokSave\%；六项任务均在评估容差内保持基线预测质量。计入全部五个截点的初始学习与记忆操作后，token 分别减少 \RbStackMavraTokSaveAll\% 与 \RbTrialMavraTokSaveAll\%。另一组三截点实验的复用阶段 token 也分别减少 \OfStackSaveReuse\% 与 \OfTrialSaveReuse\%；第~\ref{sec:relbench-eval} 节按各自配置报告两组实验。}

\bi{Our contributions are:}{本文的主要贡献如下：}
\begin{itemize}
\item \bi{A shared memory layer that turns session discoveries into executable metric and feature definitions, reducing repeated exploration and SQL construction across questions and agents.}{一个共享记忆层，将会话中的发现转化为可执行指标与特征定义，减少跨问题与智能体的重复探索和 SQL 构造。}
\item \bi{Publication, shared validation, and maintenance mechanisms that sustain reuse as data evolve and amortize checks and repairs across consumers.}{通过发布检查、共享验证与维护机制，在数据演进过程中维持复用，并跨使用者分摊检查与修复开销。}
\item \bi{A workload analysis, amortization model, and evaluation on RelBench, TPC-DS, and TPC-H connecting repeated exploration to LLM calls, token consumption, and end-to-end time, and quantifying the savings from sharing.}{通过负载分析、分摊代价模型及 RelBench、TPC-DS 和 TPC-H 上的实验，将重复探索与大模型调用、token 消耗及端到端耗时联系起来，并量化共享带来的节省。}
\end{itemize}
